\section{The first jet and marked descent}\label{sec:descent}

Let $K\subseteq k$ be the coefficient field supplied by
\cref{lem:coefficient-field}, with its chosen embedding into $\CC$.
Write $L_K,M_K$ for the resulting Lie algebras.
The marking of \cref{prop:graded-mark} is defined over $K$:
\[
 \varphi_0:\gr_FL_K\longrightarrow\gr_FM_K.
\]
Choose adapted bases whose graded classes correspond under $\varphi_0$, and
write $w_i$ for their common weights. Thus $F^a$ is spanned by the vectors
of weight at least $a$, and the marking is the identity in these graded
bases. Fix $c$ with $F^{c+1}=0$. We retain the nilradical filtration,
including $F^0g=g$, throughout this section.

\subsection{Extracting the first jet in the actual completed algebra}

Put
\[
 A=\CC[x_1,\ldots,x_n],\qquad B=A((h)),\qquad
 E=\CC((h)),\qquad R=E[[t]].
\]
An element of $B[[t]]$ has the form
\[
 \sum_{r\geq0}t^r\sum_{q\geq q_r}h^q f_{rq}(x),
 \qquad f_{rq}\in A.
\]
The lower bound $q_r$ may depend on $r$, and the coordinate degrees of
$f_{rq}$ need not be bounded. For $n>0$, this function algebra strictly
contains $E[x_1,\ldots,x_n][[t]]$: the inner series
$\sum_{q\geq0}h^q x_1^q$ belongs to $B$ but has unbounded
coordinate degree. For $n=0$, both algebras are simply $R$.
The argument below extracts coefficients of coordinate
degree zero and one before making any finite-dimensional construction;
it never assumes polynomial dependence on the coordinates over $E$.

A $t$-small vector field with coefficients in $B[[t]]$ acts as an
$R$-linear derivation of this algebra. Coordinate differentiation preserves
the lower Laurent bound; multiplication by any fixed coefficient shifts
it by a finite amount. At fixed outer order, composition involves only
finitely many outer indices. Its exponential is therefore defined
$t$-adically: at order $t^r$, only its first $r$ powers contribute.
The iterated derivation rule proves multiplicativity, and exponentiating
the negative vector field gives the inverse. These observations justify
the action on the indicated completed algebra, without requiring an
action on a larger coordinate-adic completion.

By \cref{prop:poisson-comparison}, we have an $R$-algebra isomorphism
$\Psi:B[[t]]\to B[[t]]$, congruent to the identity modulo $t$, comparing
the linear Poisson structures of the two Rees families. A common factor
$h$ in these structures can be cancelled because $h\in R^\times$.
We orient the comparison from source coordinate functions to target
coordinate functions.

\begin{lemma}\label{lem:first-jet}
The matrix
\[
 P_{ai}=\left.\partial_a\Psi(x_i)\right|_{x=0}
\]
is defined over $R$, belongs to $\GL_n(R)$, specializes to the identity,
and is a Lie isomorphism between the scalar-extended Rees families.
No preservation of augmentation by $\Psi$ is required.
\end{lemma}

\begin{proof}
Define augmentation $\epsilon:B[[t]]\to R$ and coordinate derivatives
coefficientwise in both parameters. Evaluation of each polynomial at
zero, and its coordinate derivatives, preserve its Laurent lower bound.
Consequently $\epsilon$ is an algebra homomorphism and $\partial_a$ is an
$R$-linear derivation, with
\[
 \epsilon(x_a)=0,\qquad \epsilon(\partial_a x_b)=\delta_{ab}.
\]
In particular $P_{ai}=\epsilon(\partial_a\Psi(x_i))$ is well-defined even
with unbounded coordinate degree. The marking gives $P\equiv I\pmod t$.
Writing $P=I+tS$, the geometric series $\sum_{j\geq0}(-tS)^j$ proves
invertibility over $R$ itself.

Let $C^{L,k}_{ij}(t)$ and $C^{M,r}_{ab}(t)$ denote the two Rees structure
tables. For arbitrary completed functions,
\[
 \{f,g\}_M=\sum_{a,b,r}C^{M,r}_{ab}(t)x_r
                         (\partial_a f)(\partial_b g).
\]
Its constant term vanishes. Moreover, the product rule gives
\[
 \epsilon\partial_s\{f,g\}_M
 =\sum_{a,b}C^{M,s}_{ab}(t)
                 \epsilon(\partial_a f)\epsilon(\partial_b g).
\]
Every term containing a second derivative still has the factor $x_r$
and vanishes under $\epsilon$. This finite identity establishes the
first-jet calculation directly in the completed algebra.

Apply it to the Poisson identities
\[
 \{\Psi(x_i),\Psi(x_j)\}_M
       =\sum_k C^{L,k}_{ij}(t)\Psi(x_k).
\]
If $c_i=\epsilon(\Psi(x_i))\in tR$, their constant terms give
$\sum_k C^{L,k}_{ij}(t)c_k=0$: the constant part defines a character of
the source. Their first derivatives give
\begin{equation}\label{eq:first-jet-matrix}
 \sum_{a,b}P_{ai}P_{bj}C^{M,s}_{ab}(t)
       =\sum_k C^{L,k}_{ij}(t)P_{sk}.
\end{equation}
These are exactly the Lie-matrix equations. There is no need to translate
the character to zero. Nor do we claim that $\Psi$ induces a map on the
augmentation ideal modulo its square: that ideal need not be preserved.
Invertibility came from the marking, rather than such a claim.
\end{proof}

\subsection{A regular arc with the prescribed graded special point}

Let $E_t=E((t))$. With our decreasing-filtration convention,
\[
 [x_i,x_j]_{L_t}=\sum_r c^{L,r}_{ij}
                   t^{w_r-w_i-w_j}x_r,\qquad
 \delta_{L,t}(x_i)=t^{-w_i}x_i.
\]
The latter map identifies the generic Rees algebra with
$L_K\otimes_K E_t$. The corresponding target scaling gives
\[
 f_{E_t}=\delta_{M,t}P\delta_{L,t}^{-1}:
 L_K\otimes_K E_t\xrightarrow{\sim}
 M_K\otimes_K E_t,
 \qquad (f_{E_t})_{ri}=t^{w_i-w_r}P_{ri}.
\]
Every isomorphism of the original Lie algebras preserves the nilradical
and its lower central series. Their previously established compatibility
with scalar extension therefore makes $f_{E_t}$ filtered. Hence its entries
with $w_r<w_i$ are zero. Since $R\hookrightarrow E_t$ is injective and
$t$ is invertible in $E_t$, the same entries of $P$ are zero. More explicitly, the column of weight
$w_i$ belongs to $F^{w_i}$, whose adapted basis has no vector of smaller
weight. This makes the corresponding coefficient of $f_{E_t}$ exactly
zero, rather than merely a series divisible by a power of $t$. Filtration
preservation of $P$ has thus been deduced, not imposed on the Poisson
comparison.

On equal-weight blocks the scaling factors cancel. The block diagonal
matrix
\[
 y(t)=\bigoplus_a P_{aa}(t)
\]
is a regular graded isomorphism over $R$, with $y(0)=\varphi_0$.
Here $a$ indexes weights, rather than individual basis vectors. Each block
is congruent to the identity modulo $t$, so its inverse and the inverse
of the full determinant are regular. The graded bracket equations hold
in $E_t$ and hence in $R$. The off-diagonal entries of $f_{E_t}$ may have poles;
they are never specialized at zero.

\begin{proposition}\label{prop:marked-descent}
The marked Rees comparison above yields an isomorphism
$L_K\simeq M_K$ over $K$ whose associated-graded map
is the prescribed $\varphi_0$.
\end{proposition}

We first obtain a marked point over a finite Galois extension, then prove
that its descent obstruction vanishes.

The filtered isomorphism scheme
$\mathcal I=\Iso_F(L_{\CC},M_{\CC})$ is affine of
finite type. In adapted bases it is defined by the bracket equations,
the forbidden-entry equations $f_{ri}=0$ for $w_r<w_i$, and
$z\det f=1$ with one extra coordinate $z$. An invertible block triangular
matrix has invertible diagonal blocks over every commutative coefficient
ring, since their determinant product is a unit. Thus taking diagonal
blocks defines a regular map on $\mathcal I$ and a homomorphism
\[
 \rho:\mathcal A=\Aut_F(M_{\CC})
       \longrightarrow\Aut_{\gr}(\gr_FM_{\CC}).
\]
The inverse determinant of the diagonal matrix equals $z$, so this
morphism is polynomial in the chosen affine coordinates.

The point $f_{E_t}$ makes $\mathcal I$ nonempty. The Nullstellensatz supplies
a complex point $f_*$; see \cite{StacksNullstellensatz}. Composition with
$f_*$ identifies $\mathcal I$ with $\mathcal A$. The algebraic-group
homomorphism theorem factors $\rho$ through a closed subgroup $H$, with
$\mathcal A\to H$ faithfully flat. The smooth-group version suffices
in characteristic zero by \cite[Theorems 1.72 and 3.23]{MilneAG2017}.
Each complex point of $H$ lifts, because its fiber is nonempty and of
finite type over an algebraically closed field. Consequently the graded
image of $\mathcal I(\CC)$ is the closed coset
\[
 Y=H\gr(f_*)\subseteq
 \Iso_{\gr}(\gr_FL_{\CC},\gr_FM_{\CC}).
\]
This does not assert surjectivity on rational points over $K$. In
particular, one cannot replace the following specialization argument by
the assertion that a generic isomorphism has a limit as an ordinary
isomorphism: its off-diagonal entries need not have such a limit.
Closedness is applied only after passing to the graded image, where
the entire arc, including the inverse determinant coordinate, is regular.

For completeness, the generic-point implication can be checked entirely
in finite polynomial coordinates. Let $I$ define $\mathcal I$, and let
$d$ be its diagonal-coordinate map. If a polynomial $q$ vanishes on
$Y(\CC)$, then $q\circ d$ vanishes on all complex solutions of $I$.
The Nullstellensatz gives $(q\circ d)^m\in I$ for some $m\geq1$.
Evaluation at $f_{E_t}$ implies $q(d(f_{E_t}))^m=0$, hence $q(d(f_{E_t}))=0$ in the
field $E_t$. This verifies extension-field containment without assuming
that the defining ideal is radical.

It follows that $q(y(t))\in R$ becomes zero in $E_t$, so it is already
zero in $R$. Specializing gives $q(\varphi_0)=0$ in $E$. Since $\varphi_0$ and
$q$ have complex coefficients, injectivity of $\CC\to E$ gives the same
equality over $\CC$. Thus $\varphi_0\in Y(\CC)$ and a complex filtered
isomorphism with precisely this mark exists.

The scheme $X=\Iso_{F,\varphi_0}(L_K,M_K)$ is defined
by finitely many equations over $K$ and has a complex point. Its
coordinate algebra is nonzero. A maximal ideal has residue field finite
over $K$ by \cite{StacksNullstellensatz}; enlarging that field to its
normal closure gives a finite Galois extension $K_1/K$ and a point
$f\in X(K_1)$. No algebraicity of $E$ or $E_t$ over $K$ was needed.

\subsection{The unipotent kernel and finite Galois averaging}

Let
\[
 U=\ker\bigl(\Aut_F(M_K)
               \longrightarrow\Aut_{\gr}(\gr_FM_K)\bigr).
\]
For any commutative $K$-algebra $S$, an element $u\in U(S)$ satisfies
$(u-I)F^i\subseteq F^{i+1}$ for every $i\geq0$, and therefore
$(u-I)^{c+1}=0$. The condition includes the weight-zero quotient
$M_K/N(M_K)$. Thus $U$ is a closed subgroup of a
unitriangular group, with no general linear factor left on that quotient.

Let $\mathcal A_r(S)$ be the endomorphisms raising the filtration by at
least $r$, let $U_r(S)=\{u\in U(S):u-I\in\mathcal A_r(S)\}$, and let
$D_r(S)=\operatorname{Der}_S(M_S)\cap\mathcal A_r(S)$.
For $r\geq1$ these satisfy
\[
 U_1=U,\quad U_{c+1}=1,\qquad
 \mathcal A_r\mathcal A_s\subseteq\mathcal A_{r+s}.
\]
The spaces $D_r$ are defined by linear equations over $K$,
so they commute with scalar extension. Finite polynomial formulas give
\begin{equation}\label{eq:finite-exp-log}
 \exp D=\sum_{j=0}^c\frac{D^j}{j!},\qquad
 \log u=\sum_{j=1}^c\frac{(-1)^{j+1}}j(u-I)^j.
\end{equation}
We justify, in particular, that the logarithm is a derivation. Put
$T=u-I$ and $u^s=\sum_{j=0}^c\binom{s}{j}T^j$. At nonnegative integers
$s$ this is an ordinary power of the Lie automorphism $u$, so
$u^s[a,b]=[u^sa,u^sb]$. These bounded-degree polynomials agree
identically: evaluation at sufficiently many integers has a Vandermonde
determinant which is a nonzero integer, hence a unit in every $S$.
Differentiating at zero gives the derivation rule for $\log u$, since
$\binom{s}{j}'|_0=(-1)^{j-1}/j$.

Conversely, the identity
$D^m[a,b]=\sum_j\binom mj[D^ja,D^{m-j}b]$ proves that $\exp D$
preserves brackets. Exponential and logarithm are inverse by the formal
one-variable identities evaluated on nilpotent operators. All formulas
are finite and natural in $S$. Thus $U$ is isomorphic as a scheme to the
finite-dimensional affine space $D_1$; it is smooth and
connected. Its group law need not become addition under this isomorphism.
The common finite filtration is essential here: arbitrary individually
nilpotent endomorphisms would not ensure that products of sufficiently
many different operators vanish. Here every such product raises the same
filtration beyond its last nonzero term. The inverse of $\exp D$ is
$\exp(-D)$, which also preserves the filtration and has identity graded
map. These facts establish that all corrections used below are actual
Lie automorphisms, not merely invertible linear maps.

\begin{lemma}\label{lem:unipotent-H1}
For this group $U$ over any characteristic-zero field $K$, every finite
Galois cocycle is a coboundary. In particular $H^1(K,U)=1$.
\end{lemma}

\begin{proof}
Let $\Gamma=\operatorname{Gal}(K_1/K)$ and
$a_{\sigma\tau}=a_\sigma\sigma(a_\tau)$ with $a_\sigma\in U(K_1)$.
Assume inductively that its values lie in $U_r(K_1)$. In the vector space
$\mathcal A_r/\mathcal A_{r+1}$ put
\[
 A_\sigma=[a_\sigma-I]=[\log a_\sigma].
\]
The equality holds because $2r\geq r+1$. Products of two differences
vanish in this quotient, so
$A_{\sigma\tau}=A_\sigma+\sigma A_\tau$. Define
\[
 D=\frac1{|\Gamma|}\sum_{\tau\in\Gamma}\log a_\tau
       \in D_r(K_1),\qquad v=[D].
\]
Averaging gives $\sigma v=v-A_\sigma$. For $b_r=\exp(-D)$ the transformed
cocycle $a'_\sigma=b_ra_\sigma\sigma(b_r)^{-1}$ has first-order class
$-v+A_\sigma+\sigma v=0$. Hence it takes values in $U_{r+1}$.
The transformed elements still satisfy the cocycle identity: in the
product $a'_\sigma\sigma(a'_\tau)$ the adjacent factors
$\sigma(b_r)^{-1}\sigma(b_r)$ cancel. Galois acts on the filtration
and the defining derivation equations because both are defined over $K$.
Thus the averaging and correction can be repeated with the same action.
After the finitely many stages $r=1,\ldots,c$, the transformed cocycle
is one. With $b=b_c\cdots b_1$,
then
\[
 ba_\sigma\sigma(b)^{-1}=1,\qquad a_\sigma=b^{-1}\sigma(b).
\]
Every sum is finite and its denominator is invertible in characteristic
zero. Finally, a torsor under this smooth finite-type group acquires a
point over an algebraic closure, whose finitely many coordinates belong
to a finite Galois extension. Thus the finite assertion proves the stated
Galois-cohomological conclusion.
\end{proof}

\begin{proof}[Completion of the proof of \cref{prop:marked-descent}]
The group $U$ acts on $X$ by postcomposition. For two marked points
$f_1,f_2$, their difference $f_2f_1^{-1}$ lies in $U$. These formulas
are regular and functorial, so the geometrically nonempty $X$ is a
$U$-torsor. For its point $f\in X(K_1)$ set
$a_\sigma=f(\sigma f)^{-1}$. The mark is defined over $K$, and therefore
$a_\sigma\in U(K_1)$. Direct multiplication gives
\[
 a_\sigma\sigma(a_\tau)
 =f(\sigma f)^{-1}(\sigma f)(\sigma\tau f)^{-1}
 =a_{\sigma\tau}.
\]
By \cref{lem:unipotent-H1}, write $a_\sigma=b^{-1}\sigma(b)$.
Since $\sigma f=a_\sigma^{-1}f$, we have
$\sigma(bf)=\sigma(b)a_\sigma^{-1}f=bf$.
All entries of $bf$ lie in $K$. Its determinant is nonzero, its bracket
and filtration equations hold over $K$, and its graded map is $\varphi_0$
because $b\in U$. This is the required marked isomorphism.
\end{proof}

\begin{proof}[Proof of \cref{thm:main}]
Descend the given associative equivalence, including its inverse, to the
coefficient field $K\subseteq k$ using \cref{lem:coefficient-field}.
Solvability descends to the source model. Apply the intrinsic filtration
argument over $K$, obtaining the actual marking of
\cref{prop:graded-mark}. Embed $K$ into $\CC$, construct the compatible
Rees families, and use \cref{prop:poisson-comparison,lem:first-jet} to
obtain their identity-marked Lie matrix over $\CC((h))[[t]]$.
\Cref{prop:marked-descent} gives a Lie isomorphism over $K$. Extending
scalars along $K\subseteq k$ gives $L\simeq M$ over
$k$. Thus the argument applies to every characteristic-zero ground field;
no embedding of that entire field into $\CC$ is assumed.
\end{proof}
